\section{A Unified View on Post-Training Algorithms} \label{sec:unified-view}

In this section, we adopt a unified perspective to understand both Supervised Fine-Tuning (SFT) and Reinforcement Learning (RL) as post-training objectives.
We present the gradient calculations of various post-training approaches in Table \ref{tab:grad_compare}, with exact derivations of classical approaches presented in the Appendix \ref{sec:appdendix_grad_cal}. From the table, it can be seen that policy gradient calculations for LLM post-training can be written in a unified policy gradient form. 

\begin{tcolorbox}[takeawaysbox]
We propose a unified framework for the gradient calculation of LLM post-training, named the \textcolor{darkred}{Unified Policy Gradient Estimator}:
$$ \text{grad}_{Uni} = \mathbb{1}_{stable} \frac{1}{\pi_{ref}} \hat{A} \nabla \pi_{\theta}.$$All gradient calculations can be written in the unified form.
\end{tcolorbox}


In the following sections, we further show that the differences between different gradient calculations can be broken down into four distinct components.
We theoretically derive the Unified Policy Gradient Estimator from a common objective and provide a detailed analysis of its gradient components. Based on this unified perspective, we then propose the \method (\mot) algorithm.




        


\begin{table*}[!t]
    \centering
    \caption{Theoretical unified view of various post-training algorithms.}
    \label{tab:grad_compare}
    \resizebox{\textwidth}{!}{
    \begin{tabular}{cccc}
    \toprule
         \textbf{Algorithm} & \textbf{Reference Policy} &\textbf{Advantage Estimate} & \textbf{Unified Policy Gradient Estimator} \\
        \midrule
          SFT & $\pi_{ref} = \pi_\theta$& $\hat{A}_{SFT} \equiv 1$ & $\nabla \mathcal{J}_{SFT}(\theta) = \nabla \pi_\theta (\tau ) \frac{\hat{A}_{SFT} = 1}{\pi_\theta (\tau )}$\\
        \midrule
        \multicolumn{4}{c}{\textbf{Online Reinforcement Learning Methods}} \\
        \midrule
          \makecell{PPO\\\citep{schulman2017proximal}} & $\pi_{ref} = \pi_{\theta_{old}}$&
        $\hat{A}_{PPO} = \text{GAE \citep{schulman2015high}}$ & 
        $\nabla \mathcal{J}_{PPO} = \nabla \pi_\theta (\tau) \frac{\hat{A}_{PPO} \mathbb{1}_{\text{Clip}}}{\pi_{ref} (\tau)}$\\
        \midrule
          \makecell{GRPO\\\citep{shao2024deepseekmath}} & $\pi_{ref} = \pi_{\theta_{old}}$&
        $\hat{A}_{GRPO} = \frac{R(\tau_{j}) - \text{mean} (\{R(\tau_{j})\}_{G_{on}})}{\text{std} (\{R(\tau_{j})\}_{G_{on}})}$ & 
        $\nabla \mathcal{J}_{GRPO} = \nabla \pi_\theta (\tau) \frac{\hat{A}_{GRPO} \mathbb{1}_{\text{Clip}}}{\pi_{ref} (\tau)}$\\
        \midrule
         \makecell{REINFORCE\\\citep{ahmadian2024back}} & $\pi_{ref} = \pi_{\theta}$&
        $\hat{A}_{REINFORCE} = \pm 1$ & 
        $\nabla \mathcal{J}_{REF.}(\theta) = \nabla \pi_\theta (\tau) \frac{\hat{A}_{REF.}}{\pi_\theta (\tau)}$\\
        \midrule
         \makecell{CISPO\\\citep{chen2025minimax}} & $\pi_{ref} = \pi_{\theta_{old}}$&
        $\hat{A}_{CISPO} = \hat{A}_{GRPO}$ & 
        $\nabla \mathcal{J}_{CISPO} = \nabla \pi_\theta (\tau) \frac{\hat{A}_{CISPO} \mathbb{1}_{\text{CIS-Mask}}}{\pi_{ref} (\tau)}$\\
        \midrule
         \makecell{GSPO\\\citep{zheng2025group}} & $\pi_{ref} = \pi_\theta \left(\frac{\pi_{\theta_{old}} (\tau_{i,j} |q_i ) }{\pi_{\theta} (\tau_{i,j} |q_i )}\right)^{1/|\tau_{i,j}|}  $&
        $\hat{A}_{GSPO} = \hat{A}_{GRPO}$ & 
        $\nabla \mathcal{J}_{GSPO} = \nabla \pi_\theta (\tau) \frac{\hat{A}_{GSPO} \mathbb{1}_{\text{Seq-Clip}}}{\pi_{ref} (\tau)}$\\
        \midrule
        \multicolumn{4}{c}{\textbf{Offline/Online Reinforcement Learning Methods}} \\
        \midrule
         \makecell{SRFT (Offline)\\\citep{fu2025srft}} & $\pi_{ref} \equiv 1$&
        $\hat{A}_{SRFT} = \frac{R(\tau_{j}) - \text{mean} (\{R(\tau_{j})\}_{G_{on}\cup G_{off}})}{\text{std} (\{R(\tau_{j})\}_{G_{on}\cup G_{off}})}$ & 
        $\nabla \mathcal{J}_{SRFT} = \nabla \pi_\theta (\tau) \frac{\hat{A}_{SRFT}}{\pi_{ref} (\tau) = 1}$\\
        \midrule
         \makecell{LUFFY (Offline)\\\citep{yan2025learning}} & $\pi_{ref} \equiv 1$&
        $\hat{A}_{LUFFY} = \hat{A}_{SRFT}$ & 
        $\nabla \mathcal{J}_{LUFFY} = \nabla \pi_\theta (\tau) \frac{\hat{A}_{LUFFY}}{\pi_{ref} (\tau) = 1} f_{\text{shape}}'$ \\
    \bottomrule
    \end{tabular}
    }
\end{table*}


\subsection{Components of the Unified Policy Gradient Estimator}
\label{Components of the Unified Policy Gradient Estimator}


We present the Unified Policy Gradient Estimator, our unified framework for gradient calculations. In Table \ref{tab:grad_compare}, we list a series of fundamental and well-studied post-training methods, divided into SFT and two types of RL processes. Apart from providing the closed-form policy gradients of these methods, we also present the decomposition of these methods with detailed components. It can be seen that these seemingly different methods in fact share common components and that all gradients follow our proposed unified framework.

In this paper, we divide the unified gradient into four terms: \textit{stabilization mask}, \textit{reference policy}, \textit{advantage estimate}, and \textit{likelihood gradient}. We address each of the terms below.




\paragraph{Stabilization Mask $\mathbb{1}_{stable}$} Starting from PPO \citep{schulman2017proximal}, the stabilization mask was first derived as an approximation of the TRPO Algorithm \citep{schulman2015trust}.
In practice, the PPO clipping addresses the instability issue during RL training by turning off the current gradient when the current iterate is considered unsafe. In consequent works in Table \ref{tab:grad_compare}, many have provided their modifications on the stability mask, usually motivated by empirical evaluations.



\paragraph{Reference Policy Denominator $\pi_{ref}$} The second term in our unified estimator is the reference policy on the denominator. We note that our notion of reference policy differs from the commonly used rollout policy $\pi_{\theta_{old}}$, for which we provide a discussion in Section \ref{sec:component-analysis}. This denominator denotes a token-level reweight coefficient, usually in the form of an inverse probability. There are multiple choices for this coefficient. For the case of SFT, the policy denominator uses the current policy $\pi_\theta (\tau )$. This is a result of $\Lc = -\log (\pi_\theta (\tau ))$ as the objective function. For the case of PPO-style online RL algorithms, generally, the policy denominator uses the rollout policy $\pi_{\theta_{old}} (\tau )$. Due to the unavailability of $\pi_{ref} (\tau )$ in the offline demonstration dataset, most offline RL algorithms simply assume $\pi_{ref} (\tau )= 1$ for the denominator. 


\paragraph{Advantage Estimate $\hat{A}$} 
In traditional RL, the advantage evaluates the additional benefit of taking the current action given the current state. For the context of LLMs, most of the advantage estimation is sequence-level rather than token-level, and measures the quality of the current response sequence. Similar to traditional RL literature, the post-training process seeks to maximize the likelihood of generating positive sequences with high advantage and minimize negative sequences.


\paragraph{Likelihood Gradient $\nabla \pi_\theta (\tau)$} The policy gradient term is a general term which maps gradient information from the actions to the model parameters $\theta$. It is crucial for back-propagating the objective signals to the network weights, and is kept the same across all gradient calculations. 














\subsection{Derivation of the Unified Policy Gradient Estimator}
\label{sec:shared-common-objective}


We begin from a simple and common objective shared by all post-training algorithms: improve the likelihood of positive trajectories and decrease the likelihood of negative trajectories such that the total reward in expectation $\max_{\theta} \mathcal{J} (\theta) := \mathbb{E} [r(\tau|q)]$ is maximized.
From this starting point, we theoretically derive our Unified Policy Gradient Estimator. We then show that SFT and RL objectives are not in conflict, and they can be optimized jointly within a single loss.

\paragraph{Common Objective.}
We model the post-training as a process to maximize the expected success rate while keeping the model policy closely adhering to a demonstration dataset (behavior policy) $\pi_\beta$:
\begin{equation}
\label{eq:master_obj}
\begin{aligned}
\mathcal{J}_{\mu}(\theta)
&= \E_{\tau\sim \pi_\theta(\cdot\mid q)}\!\big[r(\tau\mid q)\big]
\;-\; \mu\,\mathrm{KL}\!\big(\pi_\beta(\cdot\mid q)\,\|\,\pi_\theta(\cdot\mid q)\big),
\qquad \mu\ge 0,
\end{aligned}
\end{equation}
where $q\!\sim\!\mathcal{D}$ denotes the question from a given distribution, $\tau$ denotes a trajectory, $r$ denotes the (binary/real) score, and $\pi_\beta$ denotes behavior policy from demonstration.


\paragraph{Gradient of the Common Objective.}
Differentiating and rearranging Equation~\ref{eq:master_obj} (full derivation in Appendix~\ref{app:derivation}), we obtain
\begin{equation}
\label{eq:master_grad_pi_measure}
\begin{aligned}
\nabla_\theta \mathcal{J}_{\mu}(\theta)
&= \E_{\tau\sim \pi_\theta}\!\Big[r(\tau\mid q)\,\nabla_\theta \log \pi_\theta(\tau\mid q)\Big]
\;+\; \mu\,\E_{\tau\sim \pi_\beta}\!\big[\nabla_\theta \log \pi_\theta(\tau\mid q)\big].
\end{aligned}
\end{equation}

\paragraph{From gradient to the Unified Policy Gradient Estimator.}
Applying the measure-change identity (detailed in Appendix~\ref{app:derivation}) with the reference policy $\pi_{ref}$ which we mentioned in Section~\ref{Components of the Unified Policy Gradient Estimator} and using $\nabla \log \pi_\theta=(1/\pi_\theta)\nabla \pi_\theta$ yields the gradient:
\begin{equation}
\label{eq:upge_form}
\nabla_\theta \mathcal{J}_{\mu}(\theta)
=
\E_{\tau\sim \pi_{ref}(\cdot\mid q)}
\!\left[
\frac{1}{\pi_{ref}(\tau\mid q)}\,
\widehat{A}_{uni}(\tau,q)\,
\nabla_\theta \pi_\theta(\tau\mid q)
\right],
\end{equation}
 with the unified advantage
\begin{equation}
\label{eq:A_uni}
\widehat{A}_{uni}(\tau,q)
=
\underbrace{r(\tau\mid q)}_{\widehat{A}_{\mathrm{RL}}(\tau,q)}
\;+\;
\underbrace{\mu\,\mathbb{1}\{\pi_{ref}=\pi_\beta\}\,\frac{\pi_\beta(\tau\mid q)}{\pi_\theta(\tau\mid q)}}_{\widehat{A}_{\mathrm{SFT}}(\tau,q)}.
\end{equation}
In many RL works, the raw score $r(\tau\mid q)$ is replaced by a more structured advantage to reduce variance, provide relative credit assignment within a rollout group, and stabilize step sizes. For example, GRPO uses group-wise normalization:
\begin{equation}
\widehat{A}_{\mathrm{GRPO}}(\tau_j,q) = \frac{R(\tau_j)-\mathrm{mean}(\{R(\tau)\}_{G_{\mathrm{on}}})}{\mathrm{std}(\{R(\tau)\}_{G_{\mathrm{on}}})}.
\end{equation}

When trust-region stabilization masks, as induced by PPO clipping, are inserted multiplicatively without altering the target objective, we obtain our Unified Policy Gradient Estimator:
\begin{equation}
\label{eq:upge_masked}
\begin{aligned}
\mathrm{grad}_{uni}
&=
\E_{\tau\sim \pi_{ref}(\cdot\mid q)}\!\left[
\mathbb{1}_{stable}(\tau,q)\,
\frac{1}{\pi_{ref}(\tau\mid q)}\,
\widehat{A}_{uni}(\tau,q)\,
\nabla_\theta \pi_\theta(\tau\mid q)
\right]
\\[2pt]
&=\;
\mathbb{1}_{stable}\;
\frac{1}{\pi_{ref}}\;
\hat{A}\;
\nabla \pi_\theta.
\end{aligned}
\end{equation}
The trust-region surrogate that produces the mask is given in Appendix~\ref{app:ppo-mask-subsec}.


The gradient in \eqref{eq:master_grad_pi_measure} is the sum of two terms: (i) a reward + trust-region term sampled from $\pi_\theta$ and (ii) a data-adherence (SFT) term sampled from $\pi_\beta$. Both terms map to the same estimator via \eqref{eq:upge_form}–\eqref{eq:upge_masked} by choosing $\pi_{ref}$ accordingly (e.g., $\pi_{\theta_{old}}$ for on-policy trust-region updates and $\pi_\beta$ for SFT/offline updates). Therefore, SFT and RL optimize a single Common Objective \eqref{eq:master_obj} and can be trained jointly within one loss without intrinsic conflict.



\subsection{Gradient Component Analysis} \label{sec:component-analysis}

\begin{tcolorbox}[takeawaysbox]
\begin{enumerate}[leftmargin=1em]
    \item While all algorithms share the same Common Objective, bias-variance trade-offs still exist across current instances for different components of the unified gradient estimator.
    \item We can improve the post-training process by constructing a better and more suitable estimation of the policy gradient.
\end{enumerate}


\end{tcolorbox}


Across the wide spectrum of algorithms contained in our previous discussions and Table \ref{tab:grad_compare}, it can be inferred that the four components that construct the unified gradient estimator are motivated by different procedures in the post-training process. 
To better illustrate the relationship between the different processes with the respective components of our unified gradient, we present Figure \ref{fig:upge_overview}.



We divide the post-training process of LLMs into the four steps shown in Figure \ref{fig:upge_overview}: i) First, the LLM makes the decision on its data source, either to use data from an offline demonstration dataset, from self-generated rollout data, or a mixture of both. In this process, the policy likelihood $\pi_\theta$ of the data with respect to the current LLM is generated. ii) Given the data source used for data generation, a reference policy $\pi_{ref}$ is calculated. iii) After data collection is complete, the algorithm calculates the advantage estimation $\hat{A}$ for each token/sequence. iv) Lastly, the algorithm may choose to apply an additional masking procedure $\mathbb{1}_{stable}$ to disable the gradient calculation of various tokens, which could lead to theoretical or numerical stability issues. After these four steps, the components are collected to construct the policy gradient $\text{grad}_{Uni}$, which is used to update the LLM in the system. 
Similar to GAE presented in \citep{schulman2015high}, multiple instantiations exist to estimate the policy gradient. However, different component selections introduce various degrees of bias and variance, where a trade-off is often encountered. 
We provide the following discussion on key components of the unified gradient below.






\paragraph{Reference Policy Calculation}
Practically speaking, the reference policy denominator places a weight on each token-level update such that any token with a smaller probability, often implying more significance, is weighted more. 
SFT and REINFORCE assign weights inversely proportional to the current policy $\pi_\theta$, enforcing a bigger update when the model outputs a small probability. 
On the other hand, when the data is generated with an outdated model, algorithms such as PPO assign weights inversely proportional to the rollout policy $\pi_{\theta_{old}}$, and offline RL does not assign additional weights for tokens.



Theoretically, the reference policy is usually set given the source of the dataset and/or the rollout policy. 
For online RL methods that train purely with on-policy data, such as REINFORCE \citep{ahmadian2024back}, uses $\frac{1}{\pi_\theta}$, which produces an unbiased estimate for gradient calculation. However, these methods usually suffer from high variance.
For PPO-style online RL algorithms, the reference policy refers to the rollout policy, which is a result of importance sampling.
PPO is a numerically simplified version of TRPO \citep{schulman2015trust}. PPO makes conservative updates that effectively reduce variance. However, the important sampling ratio is in fact theoretically ill-posed and could introduce systematic bias, as discussed in GSPO \citep{zheng2025group}. GSPO has also proposed a novel calculation for $\pi_{ref}$, as shown in Table \ref{tab:grad_compare}.
On the other hand, in the offline setting, the choice for reference policy $\pi_{ref}$ is limited, since the algorithm generally has no access to the rollout policy. If we are given the assumption that the offline data evenly covers the entire state-action rollout space, then the importance sampling ratio $r(\theta) = \frac{\pi_\theta (\tau ) }{\pi_{ref} (\tau )}$ reduces to $\pi_\theta (\tau )$ by setting constant $\pi_{ref} (\tau ) = 1$. Notably, it is apparent that setting $\pi_{ref} (\tau ) = 1$ introduces much bias at the cost of numerical stability.
For the SFT case, we can consider that the domain-specific dataset is generated with respect to the expert policy $\pi^\ast$; therefore, no weighted sampling is required. 
Neither of the two approaches is entirely theoretically justified, from an RL perspective; both require a lower bound on the state-action visitation of all the possible state-action pairs \citep{kakade2003sample}, which can not be satisfied due to the severely limited datasets in practice.





Apart from the strong connection to data source and sampling polices, some studies employ a hand-crafted reweight factor within the reference policy denominator. These works \citep{yan2025learning, zhang2025policy} typically find desirable token properties and purposefully place a higher/lower weight on these desirable/undesirable tokens, respectively.

\paragraph{Choice of Stabilization Mask}
The clipping operation introduced in PPO was the first to explicitly add a stop gradient operation on LLM post-training. Clipping gradient estimation where the importance sampling strays too far from $1$ is an effective approach to address high variances. However, this aggressive clipping behavior has been criticized by some to be overly conservative: Both DAPO \citep{yu2025dapo} and CISPO \citep{chen2025minimax} stated that the classical PPO approach drops all the tokens corresponding to large model updates, and that many such tokens are in fact crucial for stabilizing entropy and facilitating scalable RL. DAPO presented a slight modification to the clipping threshold, and CISPO further extended the notion of token-wise mask, where more granular tuning was introduced to decide whether gradients from specific tokens should be dropped. The recent work of \citet{cui2025entropy} has demonstrated that many existing algorithms negatively impact the output entropy during training and introduced Clip-Cov, adding another clipping mechanism to address the entropy-collapse encountered in training. While these methods demonstrated performance enhancements in practice, they also provide additional sources of bias.

On the other hand, works such as GSPO \citep{zheng2025group} have stated that the PPO-style clipping is inherently noisy and inefficient for sample exploitation: GSPO clips a much larger fraction of tokens and yet demonstrated superior training efficiency.

In addition, post-training algorithms using offline data have chosen to purposefully remove the clipping from training, mostly guided by performance. Though setting $\pi_{ref} (\tau ) = 1$ as the policy denominator does effectively reduce the instability in gradient calculations.

\paragraph{Advantage Estimation}
There are two commonly used settings for estimating the sequence-level advantage function: the fixed advantage setting and the adaptive advantage setting. The fixed setting considers $\hat{A} = \pm 1$ given the rule-based verification, which is adapted by REINFORCE and implicitly by SFT (where all sequences are positive samples). Alternatively, recent studies have focused on using adaptive advantage estimations, performing re-centering or normalization based on the performance of the current rollout group. Notably, GRPO and its variants, such as DAPO \citep{yu2025dapo} and LUFFY \citep{yan2025learning}, use unit normalization such that the advantage estimation of the group has a unit standard deviation. Other approaches, such as Dr. GRPO \citep{liu2025understanding}, RLOO \citep{ahmadian2024back}, and REINFORCE++ \citep{hu2025reinforce++}, claim that dividing the standard deviation introduces a difficulty bias and that only recentering is adequate. 

Apart from sequence-level advantage estimate $\hat{A}_{i,j}$, recent works \citep{wang2025beyond, yang2025treerpo, sun2025ktae} have also adapted a more granular token-level advantage estimate $\hat{A}_{i,j,t}$ to a varying degree of success. 



\paragraph{A Combination of Gradient Estimators}

Although bias-variance trade-offs exist for the gradient estimator, we state that, given data distribution assumptions and sufficient data samples, all policy gradient estimators covered in our framework should result in an effective direction of improvement for the Common Objective. 
To effectively reduce the variance and bias for each policy update, we can treat instances of policy gradient as different noisy measurements of the true policy gradient, and perform a weighted average to generate a more accurate gradient estimation, similar to complementary filters \citep{marantos2015uav}.


However, the complexity of LLM RLVR introduces additional challenges. The current state of the behavior policy $\pi_\theta$ and its relationship with the respective tasks also greatly impacts the bias-variance tradeoff of each instance of the gradient estimator.
For instance, RL-zero is significantly more effective for the Qwen model series compared to LlaMA, but SFT is effective for both methods \citep{zeng2025simplerlzooinvestigatingtamingzero}; SFT $\rightarrow$ RL and RL $\rightarrow$ SFT also yield significantly different results on the same LLM \citep{fu2025srft}.
We argue that for constructing a post-training algorithm with better effectiveness and efficiency, a dynamic and adaptive mechanism is crucial to construct optimal gradient components.









\subsection{Hybrid Post-Training with Performance Feedback}
\label{sec:method}


  

  



Our unified perspective above shows that different post-training losses have the same optimization objective with different characteristics.
Inspired by this view, we propose the \method (\mot) algorithm.
We use a mixed loss $\mathcal{L} = \alpha \mathcal{L}_{\mathrm{RL}} + \beta \mathcal{L}_{\mathrm{SFT}}$, which contains the weighted on-policy RL loss $\mathcal{L}_{\mathrm{RL}}$ and SFT loss $\mathcal{L}_{\mathrm{SFT}}$, to optimize the target LLM $\pi_\theta$. The weights of the two losses ($\alpha$ and $\beta$) are determined by the real-time sampling performance of the model.

\paragraph{Performance on Single Question.} 
For any question $q$ provided to the LLM, we first obtain both a supervising trajectory $\tau^\star$ and the model's performance $P$ on the question.
Specifically, we draw $n$ on-policy trajectories $\{\tau_i\}^n_{i=1} \sim \pi_{\theta}(\cdot\mid q)$ and evaluate them with a verifier $v:\tau_i\to\{0,1\}$. This verifier is the same as the rule-based reward function and the model's performance $P$ is defined as the mean of these $n$ verification scores:
\begin{equation}
\label{eq: compute r}
v(\tau_i) = R(\tau_i) =\begin{cases}
 1 & \text{if } \tau_i \text{ contains the correct answer of } q  \\
 0 & \text{otherwise}
\end{cases} 
\end{equation}
\begin{equation}
P = \frac{1}{n}\sum_{i=1}^n v(\tau_i)
\end{equation}
Intuitively, $P$ indicates how well the current policy performs on $q$ across multiple trajectories.

\paragraph{Feedback Coefficients.} Then, we obtain the coefficients of on-policy RL loss $\alpha$ and SFT loss $\beta$ based on the performance feedback:
\begin{equation}
\alpha = f(P), \quad \beta = g(P),
\end{equation}
where the $f$ and $g$ are the specific feedback functions.
Experientially, when the model demonstrates strong capability, it is advantageous to emphasize on-policy RL to foster exploration; conversely, when the model’s competence is limited, SFT should take precedence to ensure correct guidance. Consequently, $f$ ought to be positively correlated with $P$, whereas $g$ should exhibit a negative correlation. In this paper, we employ a pair of simple yet empirically effective switch functions $f$ and $g$:
\begin{equation}
\alpha = f(P) =\begin{cases}
 1 & \text{if } P>\gamma   \\
 0 & \text{if } P\leq\gamma 
\end{cases} , \quad
\beta = g(P) =\begin{cases}
 1 & \text{if } P\leq\gamma    \\
 0 & \text{if } P>\gamma  
\end{cases}
\end{equation}
The switch gate $\gamma$ enables the model to perform SFT when its performance falls below a predefined threshold, and RL otherwise.

\paragraph{Mixed Loss.} Finally, we calculate the RL loss $\mathcal{L}_{\mathrm{RL}}$ with the already generated $n$ on-policy trajectories $\tau_i$ and SFT loss $\mathcal{L}_{\mathrm{SFT}}$ with the supervising trajectory $\tau^\star$, and we use Dr. GRPO as the on-policy RL algorithm:
\begin{equation}
\label{eq: rl loss}
\mathcal{L}_{\mathrm{RL}}
=
-\frac{1}{n}\sum_{i=1}^{n}\sum_{t=1}^{|\tau_i|}
\min\!\Big(
r_{i,t}\,A_{i,t},\;
\operatorname{clip}\!\big(r_{i,t},\,1-\epsilon,\,1+\epsilon\big)\,A_{i,t}
\Big)
\end{equation}
\begin{equation}
\label{eq: sft loss}
\mathcal{L}_{\mathrm{SFT}}=
-\frac{1}{\lvert \tau^\star \rvert}
\sum_{t=1}^{\lvert \tau^\star \rvert}
\log \pi_\theta\!\left(\tau^\star_t \,\middle|\, q, \tau^\star_{<t}\right)
\end{equation}
where $r_{i,t}=\frac{\pi_{\theta}\left(\tau_{i,t}\mid q, \tau_{i,<t}\right)}{\pi_{\theta_{old}}\left(\tau_{i,t}\mid q, \tau_{i,<t}\right)}$ is the per-token importance sampling ratio, $A_{i,t}\equiv A_i =\frac{  R(\tau_i)-\mathrm{mean}\left(\left\{\, R(\tau_i)\ \middle|\ i=1,2,\ldots,n \right\}\right)}{\mathrm{std}\left(\left\{\, R(\tau_i)\ \middle|\  i=1,2,\ldots,n \right\}\right)}$ is the advantage  and $\epsilon$ is the clip gate hyperparameter.
The mixed loss is then obtained by taking a weighted average of these two losses using performance feedback coefficients $\alpha$ and $\beta$: 
\begin{equation}
\mathcal{L} = \alpha \mathcal{L}_{\mathrm{RL}} + \beta \mathcal{L}_{\mathrm{SFT}}
\end{equation}

\begin{algorithm}[t]
\caption{The \method\ (\mot) Algorithm}
\label{alg:upt}
\SetAlgoLined
\DontPrintSemicolon
\KwIn{Pretrained LLM (policy) $\pi_\theta$; SFT dataset $\mathcal{D}_{\mathrm{SFT}}=\{(q,\tau^\star)\}$ with supervising trajectories $\tau^\star$; verifier $v$; on\mbox{-}policy samples number $n$; total training steps $T$; feedback functions $f$ and $g$; learning rate $\eta$}
\KwOut{Fine\mbox{-}tuned policy $\pi_{\theta^\ast}$.}

\For{$t = 1$ \KwTo $T$}{
  \For{$i=1$ \KwTo $n$}{
    Sample trajectory $\tau_i \sim \pi_\theta(\cdot \mid q)$\;
    Evaluate with verifier (rule-based reward): \;
    \Indp
      $v(\tau_i) \leftarrow R(\tau_i) \in \{0,1\}$\;
    \Indm
  }
  $P \leftarrow \frac{1}{n}\sum_{i=1}^{n} v(\tau_i)$\;
  $\alpha \leftarrow f(P), \quad \beta \leftarrow g(P)$\; \Hash{Performance feedback on question $q$}
  
  Compute on-policy RL loss $\mathcal{L}_{\mathrm{RL}}$ using rollouts $\{\tau_i\}$ and normalized advantages derived from $\{R(\tau_i)\}$.

  Compute SFT loss $\mathcal{L}_{\mathrm{SFT}}$ on the supervising trajectory $\tau^\star$.
  
  $\mathcal{L} \leftarrow \alpha\,\mathcal{L}_{\mathrm{RL}} + \beta\,\mathcal{L}_{\mathrm{SFT}}$\; \Hash{Mixed loss with performance feedback coefficients}

  $\theta \leftarrow \theta - \eta\,\nabla_{\theta}\mathcal{L}$\
}
\Return{$\pi_{\theta^\ast}$}\;
\end{algorithm}
